% ==========================================================================
% GENERATED by scripts/split_paper.py. Do not edit this file.
%
% Source: imported_papers/Neural_Computation__CNA/NECO-D-24-00304-Source.tex
% Part 13: Dataset
%
% Edits belong upstream, in the published paper. To change what the
% chapter includes, change imports/cna.toml and re-run
% the splitter.
% ==========================================================================
\section{Dataset} \label{sec:dataset}
The dataset used for training comprises binary images measuring ($32 \times 32$) pixels, each sample depicting a straight line going through the image center, starting and ending $2$ pixels from the image boundary  (leading to $59$ distinct lines of different angles).
The images are generated when required and may vary from one epoch to another.
During each training cycle, we randomly generate $300$ image instances (each line about $5$ times).
During evaluation, we sample each of the $59$ distinct lines exactly once.
In contrast to training, the samples used for evaluation comprise local distortions in the form of additive Gaussian nose and missing line segments (subtractive noise), and we evaluate the net fragments' capability to remove these.

To introduce additive Gaussian noise to the afferent input in stage $S2$, we probabilistically flip neurons within each channel with a probability of up to $20\%$. Since we use four feature channels in our experiments, this corresponds to a probability of $1-(1-0.2)^4=59\%$ that a neuron is flipped at any spatial location. %Below, we report the probability of noise per spatial location (where $18.5\%$ noise corresponds to flipping a neuron with a $5\%$ chance, $34.4\%$ to $10\%$ chance, $47.8\%$ corresponds to a $15\%$ chance, and $59\%$ corresponds to a $20\%$ chance).
To evaluate subtractive noise, we create discontinuous lines by deleting a line segment of $1$ to $7$ pixels in the middle.

To demonstrate the generalization capability of CNA, we generate validation samples that differ from the ones used during training. Specifically, we produce binary images measuring ($32 \times 32$) pixels depicting kinked lines by generating four random points within the image and drawing straight lines between these points.

Additionally, we generate binary images measuring ($64 \times 64$) pixels by hand that represent more complex structures.
These images include line drawings of digits $0$-$9$, characters, and objects such as a house, a star, a wheel, a clock, etc. Some of these images are shown in Figure \ref{fig:compositionality}, and an extensive list can be found at \url{https://github.com/sagerpascal/cna}.
